\section*{Author code definitions and introductory main result}
\setcounter{nttheorem}{0}
\begin{ntdefinition}[Author Definition1: Locally decodable codes (LDC)]
A code $C:\Sigma^k\to\Sigma^n$ is $(q,\delta,\varepsilon)$-locally decodable if there exists a randomized algorithm $D$, called a local decoder, that is given $i\in[k]$ as input and an oracle access to $z\in\Sigma^n$, and has the following guarantee. For every $i\in[k]$, $m\in\Sigma^k$ and $z\in\Sigma^n$ such that $\operatorname{dist}(C(m),z)\leq\delta$ it holds that $\Pr[D^z(i)\ne m_i]\leq\varepsilon$. Moreover, $D$ makes at most $q$ queries to $z$.
\end{ntdefinition}
\begin{ntdefinition}[Author Definition2: Locally correctable codes (LCC)]
A code $C\subseteq\Sigma^n$ is $(q,\delta,\varepsilon)$-locally correctable if there exists a randomized algorithm $D$, called a local corrector, that is given $j\in[n]$ as input and an oracle access to $z\in\Sigma^n$, and has the following guarantee. For every $j\in[n]$, $c\in C$ and $z\in\Sigma^n$ such that $\operatorname{dist}(c,z)\leq\delta$ it holds that $\Pr[D^z(j)\ne c_j]\leq\varepsilon$. Moreover, $D$ makes at most $q$ queries to $z$.
\end{ntdefinition}
We place $z$ in the upper script in our notation $D^z(i)$ to stress that the number of symbols read from $z$ by $D$ is of importance. The parameter $q$ is called the query complexity, and $\delta$ is the local error decoding radius, in the case LDC, and local error correction radius in the case of LCC. However, we also refer to $\delta$, somewhat inaccurately, as the local distance of the code. From here on, we do not make any explicit reference to the ``global'' distance of a code and so we refer to the local distance simply as the distance. Throughout the paper, we only consider non-adaptive LDC and LCC, defined next. Informally, these are code in which the local decoder (or corrector) samples the entries to be read before the querying step takes place. Our results only hold for non-adaptive LDCs and LCCs. For ease of discussion, throughout the introduction we ignore the error parameter $\varepsilon$. More precisely, when stating our results, every LDC or LCC (both in the hypothesis as well as in the LDC or LCC guaranteed by the theorem) has constant error.
Clearly, an LCC with a systematic encoding\footnote{An encoding from messages to codewords is called systematic if the symbols of each message are embedded in its mapped codeword.} is also an LDC and so, in particular, linear LCC induce LDC.
\begin{nttheorem}[Author introductory Theorem3: Main result]
Assume one has a non-adaptive linear $(q,\delta=\Omega(1))$-LCC with block-length $n_0$ having rate $\rho=\rho(n_0)$. Then, for every integers $\ell,c\geq1$ such that $\ell^2<c<\log n_0$, one can obtain a non-adaptive linear $(q_{\mathrm{new}},\delta_{\mathrm{new}})$-LCC with block length $n\approx n_0^\ell$, having rate $\rho_{\mathrm{new}}$, where
\begin{align*}q_{\mathrm{new}}&=(cq)^{\operatorname{poly}(\ell)},\\\delta_{\mathrm{new}}&=(cq)^{-\operatorname{poly}(\ell)},\\\rho_{\mathrm{new}}&=1-(1-\rho)^\ell-O\left(\frac{\ell^2}c\right).\end{align*}
\end{nttheorem}
